\section{Differential Chromatic Refraction}
As light passes through the Earth’s atmosphere, it is refracted toward the zenith. 
The magnitude of this refraction is wavelength-dependent, with shorter (bluer) wavelengths 
experiencing stronger refraction than longer (redder) wavelengths. This chromatic dependence 
leads to small, band-dependent shifts in the measured positions of astronomical objects, 
which in turn propagate into the modeling of the point-spread function (PSF). 
Equation~\ref{eq:dcr} defines $DCR_1$ and $DCR_2$ as the components of the differential 
chromatic refraction along the $e_1$ and $e_2$ directions, respectively, expressed as 
functions of the zenith angle $z$ and the parallactic angle $q$. To compute the DCR, 
we first determine the local sidereal time for each visit and use it to calculate 
$z$ and $q$ for every star.

\begin{equation}
\label{eq:dcr}
    \begin{aligned}
        \mathrm{DCR_1} \equiv \tan^2(z)\cos(2q) \\
        \mathrm{DCR_2} \equiv \tan^2(z)\sin(2q)
    \end{aligned}
\end{equation}
